\documentclass[11pt]{article}

%====================
% PACKAGES
%====================
\usepackage[margin=1in]{geometry}
\usepackage{iftex}
\ifPDFTeX\usepackage[T1]{fontenc}\fi
\usepackage{ebgaramond}
\usepackage{enumitem}
\usepackage{hyperref}
\usepackage{titlesec}
\usepackage{microtype}

\setlength{\parindent}{0pt}
\setlist[itemize]{noitemsep, topsep=1.5pt, leftmargin=*}
\emergencystretch=2em

%====================
% SECTION FORMATTING
%====================
\titleformat{\section}
  {\LARGE\bfseries}
  {}
  {0em}
  {}

\titlespacing*{\section}
  {0pt}
  {1.8ex}
  {0.9ex}

\titleformat{\subsection}
  {\large\bfseries}
  {}
  {0em}
  {}

\titlespacing*{\subsection}
  {0pt}
  {1.4ex}
  {0.6ex}

\begin{document}

%====================
% HEADER
%====================
\begin{center}
{\Huge\bfseries Evan Hobson}

\vspace{0.25em}

\href{mailto:evanjhobson@gmail.com}{evanjhobson@gmail.com} \;|\;
\href{https://github.com/broncomcbuck}{github.com/broncomcbuck} \;|\;
\href{https://evanhobson.com}{evanhobson.com}
\end{center}

\vspace{0.8em}

%====================
% EDUCATION
%====================
\section*{Education}

\textbf{Brigham Young University}, Provo, Utah \hfill 2021 -- Present \\
B.S. Wildlife \& Wildlands Conservation \hfill \textit{Expected Graduation: April 2027}

\vspace{0.45em}

\textbf{Focus:} Quantitative ecology, wildlife movement, spatial modeling

\vspace{0.35em}

\textbf{Selected Coursework:}
Geographic Information Systems; Wildlife Behavioral Ecology; Scientific Computing;
Linear Algebra; Computational Linear Algebra; Data Structures;
Introductory Computer Programming; Principles of Wildlife \& Fishery Management

%====================
% RESEARCH INTERESTS
%====================
\section*{Research Interests}

Movement ecology; decision-making during movement; habitat selection (RSF/SSF/iSSF);
biologging and accelerometer-based behavioral
inference; applied conservation and management on working landscapes.

%====================
% RESEARCH EXPERIENCE
%====================
\section*{Research Experience}

\subsection*{BYU Rangeland Ecology Lab \hfill Provo, UT}
\textit{Research Assistant} \hfill Sep 2025 -- Present
\begin{itemize}
\item Lead analyst for a GPS-collar study of cattle movement and grazing behavior using resource- and integrated step-selection functions (RSF/iSSF)
\item Designed and maintained reproducible spatial modeling pipelines (covariate generation, mixed-effects model fitting, cross-validation, spatial prediction) in \textbf{R} and \textbf{ArcGIS}
\item Analyzed GPS and accelerometer-linked telemetry from \textbf{64 cow-year deployments} (31 cows in 2021, 33 in 2022) on a 998-ha sagebrush steppe allotment in southeastern Idaho (5-min fixes in 2021; 7-min fixes in 2022)
\item Showed that inferred cattle response to cheatgrass depends on remote-sensing product choice by comparing three annual-grass products within identical RSF structures
\item Produced management-relevant selection maps linking space use to topography, water access, annual grass (cheatgrass) cover, perennial forage, and weather covariates to inform targeted grazing strategies
\item Conducted vegetation monitoring across two large rangeland field projects (summer 2026): the long-term Sagebrush Steppe Treatment Evaluation Project (SageSTEP) and a targeted sheep-grazing study in Idaho
\item Measured plant density, biomass, and shrub volume; canopy and basal gap; foliar cover by line-point intercept (BLM AIM); and soil-stability/erosion tests across grass, forb, and shrub functional groups
\item Awarded \textbf{BYU College Undergraduate Research Award (CURA; \$4,000)} supporting analysis, manuscript preparation, and required conference presentation
\item Re-established and advanced analysis of legacy feral horse telemetry datasets for manuscript development
\end{itemize}

\subsection*{BYU Wildlife Ecology Lab \hfill Provo, UT}
\textit{Research Assistant} \hfill Jul 2024 -- Jul 2025
\begin{itemize}
\item Conducted field research on pygmy rabbit (\textit{Brachylagus idahoensis}) ecology, including camera deployment, vegetation sampling, and habitat covariate collection
\item Assisted with mule deer GPS collar retrievals, mortality investigations, and tissue collection
\item Performed cementum annuli aging and microscopy-based laboratory workflows
\end{itemize}

%====================
% SELECTED RESEARCH PROJECTS
%====================
\section*{Selected Research Projects}

\subsection*{Cattle Resource Selection in Cheatgrass-Invaded Rangelands \hfill 2025 -- Present}
\begin{itemize}
\item Lead analyst evaluating cattle use of cheatgrass during dormant-season targeted grazing on sagebrush steppe in southeastern Idaho
\item Built three canonical RSFs with identical structure, substituting only the cheatgrass covariate, to test whether inference depends on remote-sensing product (RAP 10-m invasive annual grass, RAP 30-m annual forb-and-grass, USGS REPMS \textit{Bromus tectorum})
\item Fit binomial GLMMs (glmmTMB) with a per-cow-year random cheatgrass slope and year interaction; modeled elevation, slope, northness, distance to water, and perennial forage biomass; validated with cow-grouped cross-validation and the Boyce index
\item Found that product choice drives the inferred signal: only the species-specific product (best by AIC) showed a year-specific change from avoidance to selection, which two years of data cannot attribute to a cause
\item Manuscript with coauthors for review (lead author); target \textit{Rangeland Ecology \& Management}
\end{itemize}

\subsection*{Feral Horse Movement and Space Use in Northwestern Nevada \hfill 2025 -- Present}
\begin{itemize}
\item Analyzed legacy GPS telemetry from 11 feral horses (30-min fixes; 2 months--2 years)
\item Quantified daily distance traveled and daily space use using averaged daily minimum convex polygons (MCPs)
\item Employed MCPs as a non-traditional, derived movement metric rather than for home-range inference
\item Interpreted movement patterns relative to water availability and seasonal aridity
\item Manuscript in preparation (lead or co-lead author)
\end{itemize}

%====================
% RESEARCH OUTPUTS
%====================
\section*{Research Outputs}

\textbf{Manuscripts (Lead or Co-Lead Author)}
\begin{itemize}
\item Hobson, E., et al. \textit{Cattle resource selection in cheatgrass-invaded rangelands during dormant-season grazing}. With coauthors for review; target \textit{Rangeland Ecology \& Management}.
\item Hobson, E., et al. \textit{Seasonal and monthly variation in feral horse movement and space use in northwestern Nevada}.
\end{itemize}

%====================
% OTHER SCHOLARLY WRITING
%====================
\section*{Other Scholarly Writing \& Presentations}

\textbf{Hobson, E.}
\textit{Permission to Feel, Power to Become: How emotional engagement can enable spiritual growth}. \\
BYU Student Religion Symposium; accepted for presentation and invited for submission to symposium proceedings (non--peer reviewed).

%====================
% ADDITIONAL EXPERIENCE
%====================
\section*{Additional Experience}

\subsection*{Western North American Naturalist \hfill Provo, UT}
\textit{Assistant to Editor-in-Chief} \hfill Apr 2025 -- Present
\begin{itemize}
\item Managed dispersal of \textbf{\$10,000+} in grant funds supporting journal operations
\item Coordinated \textbf{50+ peer reviewers}, improving review turnaround and editorial workflow efficiency
\item Gained experience with peer review, editorial decision-making, and scientific publishing standards
\end{itemize}

\subsection*{BYU Mathematics I.T. Department \hfill Provo, UT}
\textit{Lead Student Computer Specialist} \hfill May 2023 -- Dec 2023
\begin{itemize}
\item Led technical support for faculty and applied mathematics labs
\item Supported Linux system administration and computational lab transitions
\end{itemize}

%====================
% TECHNICAL SKILLS
%====================
\section*{Technical Skills}

\textbf{Modeling \& Inference:}
RSF, SSF/iSSF, GLMs and mixed-effects models (glmmTMB), conditional logistic regression,
use--availability designs, cluster-robust and random-effects inference, cross-validation,
Boyce index validation, log-RSS, spatial prediction

\vspace{0.1em}

\textbf{GIS \& Spatial Data:}
ArcGIS; raster/vector workflows; covariate derivation; cartographic outputs

\vspace{0.1em}

\textbf{Programming:}
R (primary); Python; Bash; Git; exposure to C++ and Julia through coursework

\vspace{0.1em}

\textbf{Telemetry \& Field:}
GPS collars; accelerometers and behavioral-state classification; camera trapping;
rangeland vegetation monitoring: line-point intercept (BLM AIM), canopy/basal gap,
density, biomass, and shrub volume, soil-stability/erosion tests, plant functional-group ID

\vspace{0.1em}

\textbf{Scientific Writing:}
Manuscript drafting; figures; peer-review workflows

\enlargethispage{\baselineskip}
\end{document}